\documentclass[CJKmath,a4paper,10pt]{ctexart}
\usepackage{geometry}
\geometry{inner=1cm,outer=1.5cm,top=2.45cm,bottom=1cm}
\usepackage[dvipsnames,svgnames,x11names,table]{xcolor}
\usepackage{g1math-note}
\begin{document}
	\section{物理习题}
\begin{liti}	
 如图甲所示，固定粗糙斜面的倾角为$37^\circ$，与斜面平行的轻弹簧下端固定在$C$处，上端连接质量为$1\mathrm{kg}$的小滑块（视为质点），$BC$为弹簧的原长。现将滑块从$A$处由静止释放，在滑块从释放至第一次到达最低点的过程中，其加速度$a$随弹簧的形变量$x$的变化规律如图乙所示（取沿斜面向下为加速度的正方向），取$\sin37^\circ=0.6$，$\cos37^\circ=0.8$，$g=10\mathrm{m/s^2}$。根据上述过程，试求出下列各量：
 \begin{enumerate}
	\item 滑块与斜面间的动摩擦因数$\mu$；
	\item 当$x=0$时滑块的动能$E_\text{k}$；
	\item 弹簧的最大弹性势能$E_\text{p}$。
\end{enumerate}


% 题目配图（图甲：斜面+弹簧+滑块）
\begin{figure}[H]
\centering
{\everymath{\displaystyle}
\begin{tikzpicture}[>=Stealth,scale=1,every node/.style={font=\small}]
	% 斜面与地面
	\begin{scope}[xshift=-5cm]
	\coordinate (O) at (0,0);
	\coordinate (D) at (4,0);
	\coordinate (S) at (4,3);
	\draw[thick,black!80] (O) -- (D);
	\draw[thick,black!80] (O) -- (S);

	% 倾角标注
	\draw[black!70] (1.2,0) arc[start angle=0,end angle=37,radius=1.2];
	\node at (1.55,0.45) {$37^\circ$};

	% 斜面上的关键点
	\coordinate (C) at (0,0);
	\coordinate (B) at (3,2.25);
	\coordinate (A) at (4,3);

	% 弹簧与滑块
	\draw[thick,decorate,decoration={coil,segment length=4pt,amplitude=2.4pt},draw=GOneAccent!80!black] ($(C)+(0,0.2)$) -- ($(A)+(0,0.2)$);
	\begin{scope}[shift={(A)},rotate=37]
		\draw[fill=white,draw=black!80] (-0.28,0) rectangle (0.28,0.42);
	\end{scope}

	% 点与文字标注
	\fill[GOneAccent!85!black] (A) circle (1.2pt);
	\fill[GOneAccent!85!black] (B) circle (1.2pt);
	\fill[GOneAccent!85!black] (C) circle (1.2pt);
	\node[above left,rotate=37] at ($(A)+(0.2,0.1)$) {$A$};
	\node[above,rotate=37] at ($(B)+(-0.1,0.3)$) {$B$};
	\node[below left] at (C) {$C$};
	\node[right] at ($(A)+(0.35,0.1)$) {滑块};

	% BC 为原长
	\draw[|<->|,black!70] ($(B)+(-0.32,0.43)$) -- ($(C)+(-0.32,0.43)$)
		node[midway,sloped,above] {$BC$为原长};
\end{scope}
	\begin{scope}[x=10cm,y=1cm,scale=0.5,xshift=10cm,yshift=2cm]
	% 图乙：a-x图像
	
	%\begin{tikzpicture}
		\coordinate (A) at (0.2,4);
		\coordinate (B) at (0,2);
		\coordinate (C) at (0.2,0);
		\coordinate (D) at (0.6,-4);
		\coordinate (E) at(0,4);
		\coordinate (F) at(0,2);
		% 坐标轴
		\draw[->,black!80] (0,-4.5) -- (0,4.5) node[above] {$a\ (\mathrm{m\cdot s^{-2}})$};
		\draw[->,black!80] (-0.1,0) -- (0.7,0) node[right] {$x\ (\mathrm{m})$};
		% 刻度
		\foreach \y in {-4,-3,-2,-1,1,2,3,4} {
			\draw[black!65] (0.02,\y) -- (-0.02,\y) node[left] {$\y$};
		}
		\foreach \x in {0.2,0.4,0.6} {
			\draw[black!65] (\x,0.05) -- (\x,-0.05) node[below] {$\x$};
		}
		\draw[very thick,GOneAccent!85!black] (A)--(B)--(C)--(D);
		\draw[densely dotted,GOneAccent!70!black] (A) -- (E);
		\node[right,text=GOneAccent!85!black] at (A) {A};
		\node[right] at(F){B};
		\node[above]at(0.2,0){Q};
		\node[below left]at(0,0){O};
\end{scope}		
\end{tikzpicture}

	\caption{\emph{图甲为斜面弹簧装置，图乙为加速度 $a$ 随形变量 $x$ 的变化规律}}
}
\end{figure}
\tcblower
\begin{enumerate}
	\item \textbf{求滑块与斜面间的动摩擦因数$\mu$}\\
	当$x=0$时，弹簧处于原长，弹力为$0$。对滑块沿斜面方向应用牛顿第二定律：
	\[
	mg\sin37^\circ-\mu mg\cos37^\circ=ma_0.
	\]
	代入$m=1\mathrm{kg}$，$g=10\mathrm{m/s^2}$，$\sin37^\circ=0.6$，$\cos37^\circ=0.8$，$a_0=2\mathrm{m/s^2}$，得
	\[
	1\times10\times0.6-\mu\times1\times10\times0.8=1\times2,
	\qquad \mu=0.5.
	\]


\item \textbf{当$x=0$时滑块的动能$E_\text{k}$}

根据\textbf{动能定理}，合外力做功等于动能变化：$W_\text{合}=\Delta E_\text{k}$。
合外力$F_\text{合}=ma$，在$a-x$图像中，合外力做功等于$m$乘以$a-x$图线与$x$轴围成的面积（高中面积法）。

从$A$（$x=0.2\mathrm{m}$）到$B$（$x=0$），$a-x$图为梯形：
上底$a_1=2\mathrm{m/s^2}$，下底$a_2=4\mathrm{m/s^2}$，高$\Delta x=0.2\mathrm{m}$。
\[
S = \frac{(a_1+a_2)}{2} \cdot \Delta x = \frac{(2+4)}{2} \times 0.2 = 0.6\ \mathrm{m\cdot s^{-2}}
\]
合外力做功：
\[
W_\text{合}=mS=1\times0.6=0.6\mathrm{J}.
\]
因滑块从静止释放，初动能为$0$，故$E_\text{k} = W_\text{合} = 0.6\mathrm{J}$。

\item \textbf{求弹簧的最大弹性势能$E_\text{p}$}

当滑块第一次到达最低点时，速度为$0$，弹簧形变量最大，设为$x_m$。
由图乙中$x>0$部分知，滑块压缩弹簧时的加速度满足
\[
a=2-10x.
\]
\textbf{思路：} 全过程中初、末速度都为$0$，故总动能变化量为$0$。由第二问可知，物体运动到$x=0$时的动能为$0.6\mathrm{J}$，因此从$x=0$到第一次到达最低点这一段，$a-x$图像对应的有向面积应为$-0.6$。
	
从$x=0$到第一次到达最低点，由动能定理可得合外力做功等于动能变化。
在$a-x$图像中，这对应于图线与$x$轴围成的有向面积乘以质量。

其中，$x=0$到$x=0.2$之间图线在$x$轴上方，围成一个三角形，面积为
\[
S_1=\frac12\times 0.2\times 2=0.2.
\]

$x=0.2$到$x=x_m$之间图线在$x$轴下方，围成一个三角形，面积为
\[
S_2=\frac12\times (x_m-0.2)\times (10x_m-2).
\]

因为此过程末动能为$0$，初动能为$E_\text{k}=0.6\mathrm{J}$，所以合外力做功为
\[
W_\text{合}=0-E_\text{k}=-0.6\mathrm{J}.
\]

又$m=1\mathrm{kg}$，故有
\[
S_1-S_2=-0.6.
\]

即
\[
0.2-\frac12(x_m-0.2)(10x_m-2)=-0.6.
\]

化简得
\[
5x_m^2-2x_m-0.6=0.
\]

整理得
\[
x_m=0.6\mathrm{m}.
\]
（负值舍去）。

故弹簧的最大弹性势能为
\[
E_\text{p}=\frac12 kx_m^2=\frac12\times10\times(0.6)^2=1.8\mathrm{J}.
\]
\end{enumerate}
\noindent
\textbf{最终答案：}
\begin{marker}
(1) $\boxed{\mu=0.5}$；(2) $\boxed{E_\text{k}=0.6\mathrm{J}}$；(3) $\boxed{E_\text{p}=1.8\mathrm{J}}$
\end{marker}
\end{liti}
\end{document}